Num de cuenta: 323111047

\section{Introducción}

Planteamiento del problema. El reto técnico de esta actividad consiste en diseñar e implementar una aplicación de agenda telefónica dinámica utilizando el lenguaje C. El problema se centra en la transición del uso de arreglos estáticos hacia la gestión de memoria dinámica para almacenar la información de contacto, que incluye nombre, teléfono y correo electrónico del usuario. Para resolver esto, es indispensable emplear estructuras de datos tipo struct y desarrollar un menú interactivo que permita agregar y visualizar contactos en tiempo real, controlando el crecimiento del almacenamiento de forma eficiente mediante el uso de punteros durante la ejecución del programa.

\vspace{0.4cm}

Motivación. La necesidad de realizar este ejercicio surge de la importancia de crear sistemas que no estén limitados por un tamaño de memoria fijo definido desde el código fuente. En la ingeniería de software, es fundamental saber administrar los recursos del sistema para que las aplicaciones sean escalables y puedan manejar volúmenes de datos variables sin desperdiciar memoria. Dominar el uso de funciones de asignación dinámica en esta etapa es un requisito crítico para fortalecer la lógica de programación y estar preparados para los retos técnicos que implica el desarrollo del Proyecto 1 de la asignatura.

\vspace{0.4cm}


Objetivos. El propósito principal es lograr la implementación exitosa de un menú funcional que gestione las operaciones de agregar y mostrar contactos dentro de una estructura dinámica. Se busca que el alumno demuestre dominio en la reserva de espacio en memoria para arreglos de estructuras, asegurando que los datos ingresados se almacenen y desplieguen correctamente en la interfaz de línea de comandos. Finalmente, se pretende validar que el programa permita una interacción fluida y organizada, garantizando la integridad de la información desde su captura hasta su visualización final en la consola.
